\begingroup\small\singlespacing
\setlength{\tabcolsep}{3pt}
\begin{longtable}{>{\raggedright\arraybackslash}p{1.1cm}>{\raggedright\arraybackslash}p{3.2cm}>{\raggedright\arraybackslash}p{1.6cm}>{\raggedright\arraybackslash}p{3.5cm}>{\raggedright\arraybackslash}p{3.7cm}>{\raggedright\arraybackslash}p{1.6cm}}
\caption{Programas pensionarios y transferencias: versiones históricas seleccionadas}\label{tab:pensiones_historia}\\
\toprule \textbf{Inicio} & \textbf{Nombre oficial} & \textbf{Ámbito} & \textbf{Requisitos relevantes} & \textbf{Cambios históricos} & \textbf{Fuente} \\
\midrule\endfirsthead
\toprule \textbf{Inicio} & \textbf{Nombre oficial} & \textbf{Ámbito} & \textbf{Requisitos relevantes} & \textbf{Cambios históricos} & \textbf{Fuente} \\
\midrule\endhead
\bottomrule\endfoot
2001 & Apoyo Alimentario, Atención Médica y Medicamentos Gratuitos para Adultos Mayores de 70 Años Residentes en el Distrito Federal & DF & Edad: 70+. Residencia DF; confirmar requisitos del instrumento inicial & 2001, antecedente documentado. Inicio en marzo de 2001; ley publicada 18/11/2003; expansión a 68--69 en 2009 & Evalúa CDMX \citeyearpar{evalua2015} \\
2001 & Pensión Alimentaria para Personas Adultas Mayores de 68 años, residentes en la Ciudad de México & CDMX & Edad: 68+. 68+ y residencia permanente CDMX no menor de tres años; solicitud y verificación & Reglas 2017. Edad reducida desde 2009, según antecedentes; no confundir con el programa federal & Gaceta CDMX \citeyearpar{cdmx2017} \\
2003 & Programa de Atención a los Adultos Mayores & Federal & Edad: 60+. Pobreza alimentaria; localidades rurales hasta 2,500 habitantes, alta/muy alta marginación; sin apoyo similar & Reglas 2003. Denominación rural frecuente en antecedentes; mantener separado de 70 y Más. Versiones posteriores del monto pendientes & DOF \citeyearpar{rural2003} \\
2007 & Programa de Atención a los Adultos Mayores de 70 años y más en zonas rurales & Federal & Edad: 70+. Localidades hasta 2,500 habitantes; no apoyo federal similar; identidad, edad y domicilio & Reglas 2007. Inicio de 70 y Más; permite tránsito desde otros apoyos sin duplicación; no prueba identidad con el programa rural de 2003 & DOF \citeyearpar{setentamas2007} \\
2007 & Programa 70 y más & Federal & Edad: 70+. Localidades hasta 30,000 habitantes; revisar incompatibilidades de la versión & Reglas 2009. Expansión: hasta 20,000 en 2008 y hasta 30,000 en 2009; 2008 corroborado por evaluación institucional & DOF \citeyearpar{setentamas2009} \\
2007 & Programa 70 y Más & Federal & Edad: 70+. Cobertura nacional; sin ingresos por jubilación/pensión; incorporación priorizada y presupuestada & Reglas 2012. Incorpora localidades mayores de 30,000; conservar padrones anteriores y reglas del año & DOF \citeyearpar{setentamas2012} \\
2013 & Programa Pensión para Adultos Mayores & Federal & Edad: 65+. Cobertura nacional; no recibir ingreso por jubilación o pensión contributiva & Reglas 2013; denominación desde 2013. Abroga reglas 70 y Más 2012. No trasladar el umbral de 1,092 de versiones posteriores a la regla inicial de 2013 & DOF \citeyearpar{pam2013} \\
2013 & Programa Pensión para Adultos Mayores & Federal & Edad: 65+. No recibir ingreso mensual superior a 1,092 por jubilación o pensión contributiva; revisar expediente y padrón & Reglas 2014. Monto aumentado a 580; regla anual separada de 2013 & DOF \citeyearpar{pam2014} \\
2013 & Programa Pensión para Adultos Mayores & Federal & Edad: 65+. No recibir pensión contributiva superior a 1,092 mensuales; requisitos del año & Reglas 2018 antes del rediseño 2019. Última versión anual previa al rediseño; no extender criterios a 2019 & DOF \citeyearpar{pam2018} \\
2019 & Pensión para el Bienestar de las Personas Adultas Mayores & Federal & Edad: 68 general; 65 indígenas; continuidad 65--67 del padrón 2018. Residencia y documentación; reglas ya no excluyen universalmente por tener pensión contributiva & Rediseño desde 2019. Nuevo nombre y diseño. Transitorio noveno conserva personas 65--67 previamente inscritas; no asumir elegibilidad general 65 en 2019 & DOF \citeyearpar{bienestar2019} \\
2019 & Programa Pensión para el Bienestar de las Personas Adultas Mayores & Federal & Edad: 68 general; 65 indígenas/afromexicanos; continuidad padrón. Residencia mexicana; identificación y criterios diferenciados del texto inicial & 2021 antes de modificación julio. Reglas iniciales vigentes desde 01/01/2021; no representan todo el año & DOF \citeyearpar{bienestar2021a} \\
2019 & Programa Pensión para el Bienestar de las Personas Adultas Mayores & Federal & Edad: 65+. 65+ general; mexicanas por nacimiento/naturalización y domicilio en México & Modificación publicada 07/07/2021. Ampliación general a 65 y mayor monto; acuerdo entra en vigor 08/07/2021, no desde enero & DOF \citeyearpar{bienestar2021b} \\
2006 & Programa de Desarrollo Humano Oportunidades: componente para adultos mayores & Federal & Edad: 70+. Integrantes de familias beneficiarias; corresponsabilidades y restricciones del componente & Componente 2006. Componente específico, distinto de todo el ingreso familiar Oportunidades y de 70 y Más & DOF \citeyearpar{oportunidades2006} \\
2014 & PROSPERA Programa de Inclusión Social & Federal & Edad: Según componentes, no edad única. Familias en pobreza y corresponsabilidades específicas; no toda transferencia es pensión & Reglas 2015. Cambio institucional desde Oportunidades; mantener concepto de cada componente. ENASEM puede agrupar etiquetas & DOF \citeyearpar{prospera2015} \\
1994 & Programa de Apoyos Directos al Campo (PROCAMPO) & Federal & Edad: No umbral general de vejez. Productores y superficies elegibles; no inferir pensión por edad del receptor & Creación jurídica 1994; antecedentes fuera de ventana. 1994 es fecha del decreto; inicio operativo previo requiere fuente específica; no confundir subsidio productivo y pensión & DOF \citeyearpar{procampo1994} \\
2014 & PROAGRO Productivo, componente del Programa de Fomento a la Agricultura & Federal & Edad: No umbral general de vejez. Requisitos de productor, predio y componente; revisar cuotas según estrato & Ejercicio 2014. Rediseño de apoyo productivo; vinculación histórica con PROCAMPO no autoriza igualdad de montos o reglas & DOF \citeyearpar{proagro2014} \\
2019 & Programa Producción para el Bienestar & Federal & Edad: No umbral general de vejez. Pequeños/medianos productores y predios; incorporación de directorios PROAGRO, PIMAF y localidades indígenas & Lineamientos 2019. Integra antecedentes distintos; no simple cambio de nombre de un único programa & DOF \citeyearpar{produccion2019} \\
2019 & Becas para el Bienestar Benito Juárez (coordinación y programas específicos) & Federal & Edad: Elegibilidad educativa; no edad única de vejez. Distinguir programas educativos, estudiante beneficiario y hogar receptor & Decreto 2019. Reasignación desde PROSPERA y otros presupuestos; no equivalencia de todos sus componentes & DOF \citeyearpar{becas2019} \\
ND & Atención a los Adultos Mayores & Jalisco & Edad: 65+. Residencia Jalisco tres años; sin pensión de igual/superior monto; criterios económicos y presupuesto & Reglas 2014. Se documenta reducción 70$\rightarrow$65 en 2013, no inicio del programa en 2013. Separado del federal & Jalisco \citeyearpar{jalisco2014} \\
ND & Atención a los adultos mayores & Jalisco & Edad: 65+. 65+ sin ingreso por pensión contributiva; revisar requisitos completos de las reglas anuales & Ficha institucional 2018. Operación institucional corroborada en 2018; no extrapolar el monto de 2014 & Jalisco \citeyearpar{jalisco2018} \\
ND & Llego a Ti; entrega domiciliaria de despensas para adultos mayores & Mérida & Edad: Adultos mayores; mínimo exacto no documentado en informe. Beneficiarios impedidos por enfermedad para recoger apoyos; cotejar reglas originales & Segundo informe 2007--2010 (2009). Ejemplo municipal: transferencia en especie, no pensión contributiva ni ingreso monetario automáticamente & Mérida \citeyearpar{merida2009} \\
ND & Pensiones del Instituto Mexicano del Seguro Social & Federal & Edad: Por contingencia y régimen; no una edad universal. Cotización y derechos del asegurado o derivados; retiro, invalidez, riesgos y supervivencia & Regímenes coexistentes en 2001--2021. Ley publicada 1995, régimen usualmente denominado 1997; año de ley no es creación de todas las pensiones IMSS & DOF \citeyearpar{lss1995} \\
ND & Pensiones del Instituto de Seguridad y Servicios Sociales de los Trabajadores del Estado & Federal & Edad: Por modalidad y régimen; no una edad universal. Servicios/cotizaciones y derechos del trabajador o beneficiario; régimen transitorio frente a cuentas individuales & Reforma 2007 y regímenes de transición en ventana. Ley 2007 no es inicio de todas las pensiones ISSSTE; el nombre institucional no identifica régimen de cada persona & DOF \citeyearpar{issste2007} \\
\end{longtable}
\noindent\footnotesize Fuente: elaboración documental, catálogo reproducible. ND: origen no verificado, distinto de ausencia de operación. Versiones y montos son históricos; no identifican recepción individual en ENASEM.\par
\endgroup
